% !TeX root = ../../../main.tex
\section{潜空间条件输运粒子滤波与更新机制分析}\label{cw:sec:ml-updates}

在~\ref{cw:sec:ml-problem}~节的潜空间模型~\eqref{cw:eq:ml-model} 下，初始粒子取 $z_0^{+,(1..N)}\sim\mathbb P_Z$，预测步为 $z_k^{-,(i)}=z_{k-1}^{+,(i)}$，四类粒子更新做法如下：
\begin{itemize}[leftmargin=2em,topsep=3pt,itemsep=2pt,parsep=0pt]
\item \textbf{BPF 与 ETPF：}按正态窗口似然计算权重
\begin{equation}\label{cw:eq:ml-weights}
w_k^{+,(i)}\propto\exp\!\left(-\frac{1}{2\sigma^2}\bigl\|y_k-\bm H_k\mathscr T_{\Leftarrow}(z_k^{-,(i)})\bigr\|^2\right),\qquad i=1,\dots,N,
\end{equation}
分别按算法~\ref{cw:alg:bpf} 作多项重采样，或按算法~\ref{cw:alg:etpf}、\ref{cw:alg:sinkhorn-basic} 作熵正则重心变换 $Z_k^+=NZ_k^-\bm\Pi_k$。
\item \textbf{逐步在线 OTPF：}每步生成模拟观测 $y_k^{(i)}=\bm H_k\mathscr T_{\Leftarrow}(z_k^{-,(i)})+n_{y,k}^{(i)}$\cite{aljarrah2024brenier}，在线交替训练网络求解~\ref{cw:sec:ml-minimax}~节的极大极小问题~\eqref{cw:eq:ml-minimax}，令 $z_k^{+,(i)}=T_k(z_k^{-,(i)},y_k)$。
\item \textbf{静态摊销 OTPF：}由命题~\ref{cw:lem:ml-static} 与式~\eqref{cw:eq:ml-amortized-pushforward}，离线在全体训练图像与随机步数 $k$ 上训练同一个映射 $T$，在线每步直接计算
\begin{equation}\label{cw:eq:ml-amortized-update}
z_k^{+,(i)}=T\bigl(z_0^{+,(i)},\bm m_k\odot y,\bm m_k\bigr),\qquad i=1,\dots,N,
\end{equation}
并按式~\eqref{cw:eq:ml-vote} 众数读出类别（算法~\ref{cw:alg:amortized-otpf}）。
\end{itemize}

% !TeX root = ../../../main.tex
% 模板原生算法样式；离线摊销训练与在线前向更新紧凑呈现。
\begin{algorithm}[H]
\caption{静态潜空间摊销条件输运粒子滤波}\label{cw:alg:amortized-otpf}
\KwIn{训练图像 $x_{1..M}$，算子 $\mathscr T_{\Rightarrow},\ell$，掩码 $\bm m_{1..K}$，测试流 $(\bm m_k\odot y,\bm m_k)_{k=1}^K$，粒子数 $N$}
\KwOut{后验粒子 $z_k^{+,(1..N)}$ 与类别估计 $\hat c_{1:K}$}
\textbf{离线阶段：}计算 $z_j=\mathscr T_{\Rightarrow}(x_j)$，$j=1,\dots,M$\;
\While{未达训练预算}{
    随机抽取样本对、独立先验与步数 $k$，加噪构造 $\tilde y_k=(\bm m_k\odot(x+n),\bm m_k)$\;
    交替沿随机梯度上升（$f$）与多步下降（$T$）更新式~\eqref{cw:eq:ml-objective} 的 $J(f,T)$\;
}
\textbf{在线阶段：}抽取 $z_0^{+,(1..N)}\sim\mathbb P_Z$\;
\For{$k=1,\dots,K$}{
    并行计算 $z_k^{+,(i)}\gets T\bigl(z_0^{+,(i)},\bm m_k\odot y,\bm m_k\bigr)$，$i=1,\dots,N$，并按式~\eqref{cw:eq:ml-vote} 输出众数 $\hat c_k$\;
}
\end{algorithm}


\begin{table}[!htb]
\centering
\caption{四类潜空间粒子更新机制的性质对比}
\label{cw:tab:ml-updates}
\vspace{0.2em}\wuhao\setlength{\tabcolsep}{6pt}
\begin{tabular}{lccc}
\toprule[1.5pt]
更新方法 & 依赖权重 & 新粒子取值范围 & 每步在线开销 \\
\midrule[1pt]
BPF & 是 & 原有粒子子集 & $O(N)$ \\
ETPF & 是 & 原有粒子的凸包 & 每轮迭代 $O(N^2)$ \\
逐步在线 OTPF & 否 & 整个潜空间 $\mathcal S_Z$ & 一次网络训练 \\
静态摊销 OTPF & 否 & 整个潜空间 $\mathcal S_Z$ & $O(N)$ 单次前向 \\
\bottomrule[1.5pt]
\end{tabular}
\end{table}

四类方法的区别汇总于表~\ref{cw:tab:ml-updates}：
\begin{itemize}[leftmargin=2em,topsep=3pt,itemsep=2pt,parsep=0pt]
\item BPF 只能在初始粒子中复制，零过程噪声下会塌缩为单点（定理~\ref{cw:thm:bpf-impoverishment}）；ETPF 局限在旧粒子的凸包内，高维空间中权重集中时同样聚向少数粒子；两类 OTPF 不用权重，能把粒子移到 $\mathcal S_Z$ 的任意位置。
\item 逐步在线 OTPF 每步用 $N$ 个粒子在线训练，多步复合 $T_k\circ\cdots\circ T_1$ 会累积近似误差，且合成观测带有 $\mathscr T_{\Leftarrow}$ 的还原误差；静态摊销 OTPF 每步都从初始粒子一步映射，并用真实图像加噪离线训练，避开了这两项误差（一般动态模型的摊销做法可参见文献\cite{aljarrah2025aotf}，此处直接由命题~\ref{cw:lem:ml-static} 导出）。
\item 逐步在线 OTPF 每步需要训练网络，ETPF 每轮迭代需 $O(N^2)$ 计算，静态摊销 OTPF 在线只需一次 $O(N)$ 前向计算。
\end{itemize}
